\chapter{Fast targets from underneath (master inequality, version 3)}
\label{chap:under}

The camera sits on top of the airframe, looking $30$--$70^\circ$ up on a pitch hinge. The
drone flies under the target. Its camera faces backwards relative to travel, so the
nose-down pitch needed for speed swings the camera up and over: at $20$\,m/s it looks back
and up, and at $55$\,m/s it looks straight up.

\begin{lemma}[Snap limit and steady start]
\label{lem:snap}
\lean{VisionTracking.TargetParams, VisionTracking.TargetTraj.StartsSteady,
  VisionTracking.TargetTraj.accExt_second_le, VisionTracking.FIR.planAcc_second_snap}
\leanok
\uses{thm:sharp-planner}
The target's jerk changes by at most $\delta\cdot\mathrm{snap}$ per step. If the target
starts in steady flight, the planned acceleration has second differences at most
$H_0\delta^2\,\mathrm{snap}$.
\end{lemma}

\begin{theorem}[Body rates componentwise]
\label{thm:body-rate}
\lean{VisionTracking.aVec, VisionTracking.aVec_ofLp, VisionTracking.aVec_ofLp_diff,
  VisionTracking.KinPlan.rate_eq_aVec, VisionTracking.psiVec_sub_le}
\leanok
\uses{thm:lifting}
For consecutive frames $b,c$ the body rate is $\delta^{-1}\psi(a)$, where $a$ is the
antisymmetric part of $[\inner{b_i}{c_j}]$. Each component of $a$ is an inner product of a
frame \emph{difference} with frame vectors. For $\norm a,\norm{a'}\le\rho<1$ and
$G=1/\sqrt{1-\rho^2}$, $\norm{\psi(a')-\psi(a)}\le G\norm{a'-a}+(G-1)\rho$.
\end{theorem}

\begin{theorem}[Lateral-locked frame]
\label{thm:lateral}
\lean{VisionTracking.latRaw, VisionTracking.lat1, VisionTracking.lat2,
  VisionTracking.cross_lat1_lat2, VisionTracking.lat_aVec_rp, VisionTracking.lat_aVec_yaw,
  VisionTracking.lat_aVec_rp_second, VisionTracking.lat_aVec_yaw_second,
  VisionTracking.lat_steps_crude, VisionTracking.lat_up_coords}
\leanok
\uses{thm:body-rate}
With thrust direction $n$ and lateral reference $\ell$, define
$b_2=\mathrm{nrm}(-\ell+\inner\ell n\,n)$, $b_1=b_2\times n$, and $b_3=n$. This is defined at
every pitch whenever $\lvert\inner\ell n\rvert\le s_b<1$. The roll and pitch components satisfy
$a_0^2+a_1^2\le\norm{\Delta n}^2$, and the yaw component satisfies
$\lvert a_2\rvert\le(\norm{\Delta\ell}+s_b\norm{\Delta n})/\sqrt{1-s_b^2}$. Their changes are
bounded by $\norm{n'\times\Delta^2n}$, $\norm{\Delta^2\ell}$ and products of first
differences.
\end{theorem}

\begin{lemma}[Target and plan at the current speed]
\label{lem:local}
\lean{VisionTracking.norm_targetAcc_le_loc, VisionTracking.thrustStar_up_loc,
  VisionTracking.thrustStar_along_v_ge, VisionTracking.thrustStar_along_heading_ge,
  VisionTracking.abs_inner_targetAcc_perp_le, VisionTracking.Under.thr_err,
  VisionTracking.Under.thr_norm_ge, VisionTracking.Under.lateral_le,
  VisionTracking.Under.forward_ge, VisionTracking.Under.n_cross_second,
  VisionTracking.Under.h_second, VisionTracking.Under.window}
\leanok
\uses{lem:snap}
All bounds take the current target speed $s$ and hold over the three steps one control step
uses. At high speed the drag makes the thrust large, so the bank $\arcsin s_b$ stays near
$30^\circ$ even in a $350$\,m turn at $55$\,m/s.
\end{lemma}

\begin{lemma}[Hinge geometry]
\label{lem:hinge}
\lean{VisionTracking.clamp_lip, VisionTracking.clamp_dist, VisionTracking.inner_cam_up_eq,
  VisionTracking.abs_arctan_sub_le}
\leanok
With in-plane pitch $\varphi=\arctan(u_1/u_3)$, a hinge angle $\eta$ gives
$\inner{\mathrm{cam}}{e_3}=\rho\cos(\pi/2-\eta-\varphi)$. The clamp to the hinge range is
1-Lipschitz, and $\varphi$ moves by at most $\tfrac\pi2\norm{\Delta(u_3,u_1)}/\rho_{\min}$.
\end{lemma}

\begin{definition}[Master budget, version 3]
\label{def:budget3}
\lean{VisionTracking.MasterIneq3, VisionTracking.Budget3.sb, VisionTracking.Budget3.θa,
  VisionTracking.Budget3.yaw2, VisionTracking.Budget3.τrp, VisionTracking.Budget3.τyaw,
  VisionTracking.Budget3.camOff, VisionTracking.Budget3.hingeStep}
\leanok
\uses{thm:lateral, lem:local, lem:hinge}
Conditions at every speed of the band: positive vertical thrust, a bank below $90^\circ$,
close consecutive frames, tilt, body rate, mixer (low and high), hinge rate, and the
visibility condition $\mathrm{camOff}+\mu\le\alpha$, with
$\mathrm{camOff}=\arccos(\sqrt{1-s_b^2}\cos\mathrm{miss})$.
\end{definition}

\begin{theorem}[Master inequality, version 3, sufficient side]
\label{thm:master3}
\lean{VisionTracking.master_sufficient3, VisionTracking.MaintainableSteady,
  VisionTracking.SteadyStrategy, VisionTracking.Under.traj, VisionTracking.Under.traj_keeps,
  VisionTracking.Under.traj_causal, VisionTracking.Under.kp, VisionTracking.Under.kp_mix,
  VisionTracking.Under.dav_le, VisionTracking.Under.B_step_le, VisionTracking.Under.eta_rate,
  VisionTracking.Under.cam_up_angle}
\leanok
\uses{def:budget3, thm:master2}
If \texttt{MasterIneq3} holds for a planner with four frames of lag, a position-only
strategy keeps every target that starts in steady flight in view in every frame. It flies
underneath as a valid rigid-body quadrotor.
\end{theorem}

\begin{remark}[Where it holds]
\texttt{scripts/budget\_v3.py} mirrors \texttt{Budget3} definition by definition. The
setup is a $20$--$55$\,m/s target with turn radius $350$\,m, snap $15$\,m/s$^4$ and
vertical acceleration $1.5$\,m/s$^2$; a $1.5$\,kg quadrotor; and a sparse three-tap
quadratic-exact planner.
\begin{itemize}
  \item With the base quadrotor, \texttt{MasterIneq3} holds for jerk at most
        $3$\,m/s$^3$ with $5^\circ$ dives, and jerk $1$ with $10^\circ$ dives.
  \item With yaw coefficient $0.03$ and a $22$\,cm arm, it holds for jerk at most $3$ and
        $2$ respectively.
  \item Jerk $5$ fails on torque at low collective thrust.
\end{itemize}
\end{remark}

\chapter{Looking down from above}
\label{chap:lookdown}

\begin{lemma}[Yaw-free look-down tilt]
\label{lem:lookdown}
\lean{VisionTracking.QuadTraj.angle_b₃_cam, VisionTracking.QuadTraj.tilt_ge_of_visible,
  VisionTracking.angle_up_ge_of_below, VisionTracking.QuadTraj.tilt_ge_of_visible_below}
\leanok
A target at least $h$ below the drone, within detection range $D$, forces the tilt to be at
least $\arcsin(h/D)+\eta_{\mathrm{lo}}-\alpha$, whatever the yaw.
\end{lemma}

\begin{theorem}[Free tilt from drag]
\label{thm:lookdown-steady}
\lean{VisionTracking.steady_look_down_iff}
\leanok
\uses{lem:lookdown}
Behind a straight level target at speed $v$, a depression $e$ is held forever iff
$e\le\mathrm{chaseTilt}(v)+\alpha-\eta_{\mathrm{lo}}$ (together with the hinge's upper side).
\end{theorem}

\begin{theorem}[Forced loss]
\label{thm:forced-loss}
\lean{VisionTracking.QuadTraj.forced_loss_time, VisionTracking.planar_tilt,
  VisionTracking.angle_ge_of_opposite, VisionTracking.QuadTraj.angle_b₃_step_le,
  VisionTracking.inner_eq_track_coords}
\leanok
\uses{lem:lookdown}
Suppose a drone keeps its thrust axis in the vertical plane of the track, cannot flip its
tilt in one step, and keeps its velocity within $U$ of the target's. If it must keep a
target at least $h$ below it in view, then it does so for at most
$2U(\cos p_0+\sin p_0)/\kappa$, with $\kappa=\sin p_0(g-cVU)-cV^2\cos p_0$. Orbiting or
weaving drones are not covered.
\end{theorem}

\chapter{P3: the necessary side}
\label{chap:p3}

\begin{theorem}[Directional speed wall, any trajectory]
\label{thm:dirspeed-wall}
\lean{VisionTracking.no_traj_keeps_of_dirSpeed_lt, VisionTracking.straightTarget_startsSteady,
  VisionTracking.not_maintainableSteady_of_dirSpeed_lt}
\leanok
A target flying straight along an admissible direction $u$ faster than the drone's
drag-limited directional speed $\mathrm{dirSpeed}\,u$ escapes \emph{every} drone trajectory,
even one planned with knowledge of the target's whole future. The escaping target starts in
steady flight.
\end{theorem}

\begin{theorem}[P3 bracket]
\label{thm:p3-bracket}
\lean{VisionTracking.NecIneq3, VisionTracking.master_necessary3, VisionTracking.p3_bracket}
\leanok
\uses{thm:master3, thm:dirspeed-wall}
$\texttt{MasterIneq3}\Rightarrow$ every steadily starting target can be kept in view
$\Rightarrow\texttt{NecIneq3}$, where $\texttt{NecIneq3}$ requires
$v_{\max}\le\mathrm{dirSpeed}\,u$ in every admissible direction.
\end{theorem}

\begin{remark}[What does not help, and turns]
Lag-compensating linear planners cannot reduce the torque the plan needs. With four
frames of lag they must extrapolate, which forces a gain $\sum|h_i|\gtrsim2$, so smoothing
never beats the target's own jerk and snap limits. A sustained-turn thrust necessity is
vacuous here: the detection range (about $570$\,m) exceeds the turn radius ($350$\,m), so
a drone hovering near the turn centre keeps a circling target in view.
\end{remark}

\begin{remark}[The gap]
\texttt{scripts/gap\_v3.py} compares the two sides. For the reference quadrotor, the
necessary side allows targets up to about $85$\,m/s. The sufficient side reaches about
$64$\,m/s with $5^\circ$ dives: for jerk at most $5$ with yaw coefficient $0.03$ and a
$22$\,cm arm, and for jerk at most $3$ with the base quadrotor.
The torque part of \texttt{MasterIneq3} drives the gap. It reflects the design choice of
copying the target's acceleration exactly; it is not a physical limit.
\end{remark}
